\section{Local divisibility with its exact lattice costs}\label{sec:local-accounting}

\status{The deductions below are proved here. The joint exceptional-rank estimates and the valuations of the actual theta operators remain arithmetic inputs. The index estimate in the source's local selection theorem now has an explicit elementary-divisor proof.}

\subsection{From exceptional directions to a lattice index}

Let $\mathcal O$ be a discrete valuation ring with uniformizer $\pi$. Let $M$ be free of rank $r>0$, and let $N\subseteq M$ be a full lattice. For $a\ne0$ in the maximal ideal put $e=v_\pi(a)$. Say that $N$ has \emph{at most $g$ exceptional directions at depth $E$} if there are $g$ vectors in $M$ whose span, together with $N$, contains $a^E M$. This is the source's \texttt{AllButLE}; using fewer than $g$ vectors makes no difference.\cite{local-selection,local-smith}

Choose a Smith basis, so $M=\bigoplus_i\mathcal O e_i$ and $N=\bigoplus_i\pi^{c_i}\mathcal O e_i$. The condition is equivalent to
\[
\#\{i:c_i>eE\}\le g.
\]
For one direction, retain the corresponding basis vectors. Conversely, the summands with $c_i>eE$ require that many generators after reduction at depth $eE+1$; a submodule of a $g$-generated module over a DVR has at most $g$ generators. This is the elementary-divisor content of the formal statement. \textsf{[LC1]}

\begin{proposition}\label{prop:local-index}
Suppose $\tau,S\ge0$ are integers and $\kappa_v\ge0$ is a finite family. Write
\[
e_j(\kappa)=\sum_{|B|=j}\prod_{v\in B}\kappa_v,
\qquad \Sigma=m\sum_v\kappa_v,
\]
where $m\ge1$ is an integer. Suppose that for every $\beta\ge0$ the number of exceptional directions at depth $\tau+\beta S$ is at most
$r m^{\beta+1}e_{\beta+1}(\kappa)$. Then, for $N_*:=N\cap a^\tau M$,
\begin{equation}\label{eq:local-index-exact}
\operatorname{length}_{\mathcal O}(M/N_*)
\le er\left[\tau+S\left\{\prod_v(1+m\kappa_v)-1\right\}\right]
\le er\{\tau+S(e^\Sigma-1)\}.
\end{equation}
In particular, if $\Sigma\le1/2$, then the length is at most $er(\tau+2S\Sigma)$.
\end{proposition}

\begin{proof}
The elementary divisors of $N_*$ have exponents $\max(c_i,e\tau)$, so the length is
\[
er\tau+\sum_i(c_i-e\tau)_+.
\]
For $S>0$, divide the tail of each exponent into intervals of length $eS$. On the $\beta$th interval, the number of surviving exponents is at most $r m^{\beta+1}e_{\beta+1}(\kappa)$. Thus the sum of the tails is at most
\[
eSr\sum_{j\ge1}m^j e_j(\kappa)
=eSr\left\{\prod_v(1+m\kappa_v)-1\right\}.
\]
The sums are finite, since $e_j(\kappa)=0$ beyond the number of indices. If $S=0$, this eventual zero bound forces every $c_i\le e\tau$, giving the same conclusion. Finally $1+y\le e^y$, and for $0\le\Sigma\le1/2$,
$e^\Sigma-1\le\Sigma/(1-\Sigma)\le2\Sigma$.
\end{proof}

The last inequality is the source's index allowance. The product and exponential bounds in \eqref{eq:local-index-exact} retain information that its proof discards; these are separately derived improvements. There is also a useful saturated refinement. If $M'\subseteq M$ is saturated of rank $r'$, then
\[
\operatorname{length}\bigl(M'/(M'\cap N_*)\bigr)
\le er'\tau+erS\left\{\prod_v(1+m\kappa_v)-1\right\}.
\]
The baseline $er'\tau$ is exact; the remaining quotient embeds in $a^\tau M/(N\cap a^\tau M)$, whose length is the tail sum just estimated. Applying this refinement to the final automorphic lattices requires checking their reference-lattice transports; we do not silently replace the source's established allowance. \textsf{[LC2]}

In the actual local selection argument, the candidate lattice is generated by
\[
a^{\sum_{v\notin B}S_v}\Bigl(\prod_{v\in B}y_v^m\Bigr)M,
\qquad B\subseteq\mathcal V,
\]
where the $y_v$ commute, $S_v\le S$, and
$\tau=\sum_{v\in\mathcal V}m(t_v+D_0+1)$. If necessary replace the generated lattice by its intersection with $M$; this leaves both the exceptional-direction hypothesis and $N_*$ unchanged. The source proves that hypothesis using its rational-plane estimates and compatible lattice systems. Proposition~\ref{prop:local-index} does \emph{not} prove those estimates. It explains exactly how they pay for restricting to $N_*$.\cite{local-systems}

\subsection{Normalization and rounding over a number field}

Let $K$ be a number field of degree $k_0$, and for a finite place $w$ write
\[
\ell_w=\frac{\log Nw}{k_0},\qquad e_w=v_w(p)\quad(w\mid p),
\qquad \operatorname{sz}_w(y)=2\ell_w e_w y.
\]
The fundamental identity $\sum_{w\mid p}e_w f_w=k_0$ gives
\begin{equation}\label{eq:prime-normalization}
\sum_{w\mid p}\operatorname{sz}_w(y)=2y\log p.
\end{equation}
No Galois hypothesis is needed for this identity. It explains how local normalized gains become rational-prime sums without an extra factor from the degree of the coefficient field. \textsf{[LC3]}\cite{local-normalization}

For the slope statements take $K$ totally complex and use the source's squared-norm convention. For finitely supported real shifts $s_u(w)$, set
\[
I_u=\prod_w\mathfrak p_w^{-\lfloor e_ws_u(w)\rfloor},\qquad
\mathcal R_u=2\sum_w\ell_w\{e_ws_u(w)\},
\]
where $\{z\}=z-\lfloor z\rfloor\in[0,1)$, also for negative $z$. The slope identity, with the same archimedean metrics, is exact:
\begin{equation}\label{eq:round-exact}
\mu(I_uN_u)-\mu(N_u)
=\sum_w\operatorname{sz}_w(s_u(w))-\mathcal R_u.
\end{equation}
Indeed multiplication by $I_u$ multiplies covolume by $N(I_u)^{\operatorname{rank}N_u}$; take its normalized logarithm. Thus rounding is a nonnegative loss, not an unsigned unspecified error. The floor is taken \emph{after} multiplying by $e_w$.\cite{local-div-defs}

If the nonzero rounding terms lie over primes $p\le X$ and at most three exceptional primes $g_1,g_2,T_u$, then
\[
\mathcal R_u\le2\left(\vartheta(X)+\log g_1+\log g_2+\log T_u\right),
\qquad \vartheta(X)=\sum_{p\le X}\log p.
\]
Here $\sum_{w\mid p}\ell_w\le\log p$, since $e_w\ge1$. The elementary bound $\vartheta(X)\le X\log4$ follows by induction from the central binomial coefficients: use $\binom{2n}{n}\le2^{2n}$ for even arguments and $\binom{2n+1}{n}\le2^{2n}$ for odd ones, with the prime intervals $(n,2n]$ and $(n+1,2n+1]$, respectively. Therefore if each exceptional logarithm is at most $X$,
\begin{equation}\label{eq:round-source}
0\le\mathcal R_u\le(2\log4+6)X.
\end{equation}
This reconstructs the source's explicit rounding allowance. Its required support and type-prime bounds remain hypotheses until supplied by the chosen family. \textsf{[LC4]}\cite{local-rounding}

\subsection{The precise net-gain identity}

The source defines the real shift by
\[
s_u(w)=\sum_j\{g_j(w)+F_j(w)-E_{u,j}(w)\},
\]
where $g_j$ already incorporates the discarded-prime convention. Define
\begin{align*}
\mathrm{Ga}&=\sum_w\operatorname{sz}_w\!\left(\sum_jg_j(w)-\tau_w\right),&
\mathrm{Fo}&=\sum_w\operatorname{sz}_w\!\left(\sum_jF_j(w)\right),\\
\mathrm{Ee}_u&=\sum_w\operatorname{sz}_w\!\left(\sum_jE_{u,j}(w)\right),&
\mathrm{IO}&=\sum_w\operatorname{sz}_w(\mathrm{io}_w).
\end{align*}
Also put $\mathrm{Ta}=\sum_w\operatorname{sz}_w(\tau_w)$. Let $M_u$ be the retained reference source lattice of rank $r_u>0$, and let $d_u\ge r_u$ be the full column rank. Put $\rho_u=d_u/r_u$. The source bounds the local restriction index using the \emph{full} rank:
\[
\operatorname{length}((M_u)_w/(N_u)_w)
\le e_w d_u(\tau_w+\mathrm{io}_w).
\]
Assume this at the regular places, and write $\mathcal B_u$ for the remaining nonnegative bound on the excess signed relative index at exceptional places. The determinant formula for lattice index, followed by \eqref{eq:round-exact}, gives
\begin{equation}\label{eq:net-gain-exact}
\begin{split}
\mu(I_uN_u)\ge{}&\mu(M_u)+\mathrm{Ga}+\mathrm{Fo}-\mathrm{Ee}_u\\
&-\rho_u\mathrm{IO}-(\rho_u-1)\mathrm{Ta}-\mathcal R_u-\mathcal B_u.
\end{split}
\end{equation}
This displays the dimension ratio that cannot be dropped when passing from full columns to retained columns. If $r_u\ge(1-d)d_u$, $0\le d<1$, and $\tau_w,\mathrm{io}_w\ge0$, the two index terms are bounded by
\[
\rho_u\mathrm{IO}+(\rho_u-1)\mathrm{Ta}
\le\frac{\mathrm{IO}}{1-d}+\frac{d\,\mathrm{Ta}}{1-d}.
\]
This is exactly the structure of the source's \texttt{ClaimLow}. At an exceptional place use the signed relative index
\[
\operatorname{length}(M_u/(M_u\cap N_u))_w
-\operatorname{length}(N_u/(M_u\cap N_u))_w;
\]
its nonnegative excess beyond $e_wd_u(\tau_w+\mathrm{io}_w)$, divided by $r_u$ and weighted by $2\ell_w$, defines $\mathcal B_u$. This convention also covers a local enlargement instead of a restriction. \textsf{[LC5]}

In the released construction the word's own type prime $T_u$ is explicitly excluded from the regular-place restriction hypothesis. Its contribution must therefore be bounded separately. The source packages this and other permitted losses in \texttt{Lo5}; the present argument does not set them to zero. Equations \eqref{eq:round-exact}--\eqref{eq:net-gain-exact} give the precise bookkeeping behind the schematic net gain of version 0.1.\cite{local-div-defs,local-costs}

\paragraph{Absorbing the ramified rank and mean corrections.}
For an integer $f\ge0$, the source also supplies the separate identities
\[
\frac{r_u}{d_u}=\left(1-\frac1{T_u-1}\right)^f,
\qquad
\mu(M_u)=\mu(M_u^{\rm full})-
\frac{f\log(T_u-1)}{T_u-2},\qquad T_u>2.
\]
The first gives $\delta_u:=1-r_u/d_u\le f/(T_u-1)\le2f/T_u$ by Bernoulli's inequality. The second is an additional metric input, not a consequence of the rank identity. Assume $\mathrm{IO}\ge0$, $0\le\mathrm{Ta}\le b^2$, $2f/T_u\le1/2$, and
\[
\frac{4fb^2}{T_u}+\frac{f\log(T_u-1)}{T_u-2}\le1.
\]
Then $\mathrm{IO}/(1-\delta_u)\le2\mathrm{IO}$, while
$\delta_u\mathrm{Ta}/(1-\delta_u)\le4fb^2/T_u$.
Thus \eqref{eq:net-gain-exact} implies
\[
\mu(I_uN_u)\ge\mu(M_u^{\rm full})+\mathrm{Ga}+\mathrm{Fo}-\mathrm{Ee}_u
-\bigl(2\mathrm{IO}+\mathcal B_u+\mathcal R_u+1\bigr).
\]
When $\mathrm{io}_w$ depends only on the rational prime $p$ below $w$, put $\mathrm{IOp}=\sum_p\mathrm{io}_p\log p$. Equation~\eqref{eq:prime-normalization} gives $\mathrm{IO}=2\mathrm{IOp}$, yielding the numerical $4\mathrm{IOp}$ allowance in \texttt{ClaimLow2}. This accounts for the ramified corrections explicitly without asserting their underlying metric theorem. \textsf{[LC6]}
